% ======================================================================
% Beamer presentation - generated by KherveSlide.
% Edit freely; the source is regenerated when you change a slide.
% ======================================================================
\documentclass[aspectratio=169]{beamer}

% --- theme & packages ---
\usetheme{default}
\setbeamertemplate{itemize items}[ball]
\definecolor{ksth0}{HTML}{002147}
\setbeamercolor{structure}{fg=ksth0}
\definecolor{ksth1}{HTML}{212121}
\setbeamercolor{normal text}{fg=ksth1}
\definecolor{ksth2}{HTML}{002147}
\setbeamercolor{frametitle}{fg=ksth2}
\definecolor{ksth3}{HTML}{002147}
\setbeamercolor{title}{fg=ksth3}
\definecolor{ksth4}{HTML}{D1D7DE}
\setbeamercolor{block title}{bg=ksth4}
\definecolor{ksRule}{HTML}{A79D96}
\addtobeamertemplate{frametitle}{}{%
\vskip2pt{\color{ksRule}\hrule height 1.5pt}}
\usepackage{lmodern}
\usepackage{tikz}
\usetikzlibrary{arrows.meta}
\usepackage[absolute,overlay]{textpos}
\usepackage{graphicx}
\setlength{\TPHorizModule}{1\paperwidth}
\setlength{\TPVertModule}{1\paperheight}
\textblockorigin{0\paperwidth}{0\paperheight}
\setbeamerfont{itemize/enumerate subbody}{size=\relax}
\setbeamerfont{itemize/enumerate subsubbody}{size=\relax}
\title{The Gaussian integral}
\author{Analysis seminar}

\begin{document}

% ======================================================================
% Slide 1
% ======================================================================
\begin{frame}[t]

  % text box (free-positioned)
\begin{textblock}{0.84}(0.08,0.16)
{\centering\fontsize{32}{38}\selectfont \textcolor[HTML]{1A1A1A}{\textbf{The Gaussian integral}}\par}
\end{textblock}

  % line / arrow (free-positioned)
\definecolor{ksline0}{HTML}{002147}
\begin{textblock}{0.24}(0.38,0.52)
\begin{tikzpicture}
\useasboundingbox (0,0) rectangle (\linewidth,-0\TPVertModule);
\draw[line width=2.5pt,color=ksline0] (0\linewidth,-0\TPVertModule) -- (1\linewidth,-0\TPVertModule);
\end{tikzpicture}
\end{textblock}

  % text box (free-positioned)
\begin{textblock}{0.84}(0.08,0.57)
{\centering\fontsize{17}{20}\selectfont \textcolor[HTML]{555555}{Three proofs and a picture}\par}
\end{textblock}

  % text box (free-positioned)
\begin{textblock}{0.84}(0.08,0.76)
{\centering\fontsize{13}{16}\selectfont \textcolor[HTML]{555555}{Analysis seminar}\par}
\end{textblock}
\end{frame}

% ======================================================================
% Slide 2 - Statement
% ======================================================================
\begin{frame}[t]
  \frametitle{Statement}

  % text box (free-positioned)
\begin{textblock}{0.88}(0.06,0.22)
{\begin{definition}\raggedright\fontsize{15}{18}\selectfont $f(x) = e^{-x^2}$ is integrable on $\mathbb{R}$.\end{definition}\par}
\end{textblock}

  % text box (free-positioned)
\begin{textblock}{0.88}(0.06,0.4)
{\begin{theorem}[Gauss]\raggedright\fontsize{18}{22}\selectfont \[\int_{-\infty}^{\infty} e^{-x^2}\,dx = \sqrt{\pi}\]\end{theorem}\par}
\end{textblock}

  % text box (free-positioned)
\begin{textblock}{0.88}(0.06,0.74)
{\begin{proof}\raggedright\fontsize{15}{18}\selectfont Square it, then switch to polar coordinates.\end{proof}\par}
\end{textblock}
\end{frame}

% ======================================================================
% Slide 3 - Proof by polar coordinates
% ======================================================================
\begin{frame}[t]
  \frametitle{Proof by polar coordinates}

  % text box (free-positioned)
\begin{textblock}{0.66}(0.04,0.2)
{\centering\fontsize{22}{26}\selectfont \[I^2 = \int\int e^{-(x^2+y^2)}\,dx\,dy\]\par}
\end{textblock}

  % text box (free-positioned)
\begin{textblock}{0.22}(0.74,0.24)
{\raggedright\fontsize{15}{18}\selectfont \textcolor[HTML]{666666}{\textit{square}}\par}
\end{textblock}

  % text box (free-positioned)
\begin{textblock}{0.66}(0.04,0.39)
{\centering\fontsize{22}{26}\selectfont \[= \int_0^{2\pi}\int_0^{\infty} e^{-r^2}\,r\,dr\,d\theta\]\par}
\end{textblock}

  % text box (free-positioned)
\begin{textblock}{0.22}(0.74,0.43)
{\raggedright\fontsize{15}{18}\selectfont \textcolor[HTML]{666666}{\textit{polar}}\par}
\end{textblock}

  % text box (free-positioned)
\begin{textblock}{0.66}(0.04,0.58)
{\centering\fontsize{22}{26}\selectfont \[= 2\pi \cdot \frac{1}{2} = \pi\]\par}
\end{textblock}

  % text box (free-positioned)
\begin{textblock}{0.22}(0.74,0.62)
{\raggedright\fontsize{15}{18}\selectfont \textcolor[HTML]{666666}{\textit{integrate}}\par}
\end{textblock}

  % text box (free-positioned)
\begin{textblock}{0.9}(0.04,0.84)
{\centering\fontsize{18}{22}\selectfont \textcolor[HTML]{002147}{\textbf{Hence $I = \sqrt{\pi}$.}}\par}
\end{textblock}
\end{frame}

% ======================================================================
% Slide 4 - The picture behind it
% ======================================================================
\begin{frame}[t]
  \frametitle{The picture behind it}

  % shape (free-positioned)
\definecolor{ksshape1}{HTML}{002147}%
\begin{textblock}{0.32}(0.1,0.22)
\begin{tikzpicture}
\useasboundingbox (0,0) rectangle (\linewidth,-0.57\TPVertModule);
\path[draw=ksshape1, line width=1.2pt] (0.5\linewidth,-0.285\TPVertModule) ellipse [x radius=0.5\linewidth, y radius=0.285\TPVertModule];
\end{tikzpicture}
\end{textblock}

  % line / arrow (free-positioned)
\definecolor{ksline2}{HTML}{888888}
\begin{textblock}{0.4}(0.06,0.505)
\begin{tikzpicture}
\useasboundingbox (0,0) rectangle (\linewidth,-0\TPVertModule);
\draw[line width=1pt,color=ksline2, -{Stealth[length=2.4mm]}] (0\linewidth,-0\TPVertModule) -- (1\linewidth,-0\TPVertModule);
\end{tikzpicture}
\end{textblock}

  % line / arrow (free-positioned)
\definecolor{ksline3}{HTML}{888888}
\begin{textblock}{0}(0.26,0.18)
\begin{tikzpicture}
\useasboundingbox (0,0) rectangle (\linewidth,-0.66\TPVertModule);
\draw[line width=1pt,color=ksline3, -{Stealth[length=2.4mm]}] (0\linewidth,-0.66\TPVertModule) -- (0\linewidth,-0\TPVertModule);
\end{tikzpicture}
\end{textblock}

  % line / arrow (free-positioned)
\definecolor{ksline4}{HTML}{C8102E}
\begin{textblock}{0.12}(0.26,0.305)
\begin{tikzpicture}
\useasboundingbox (0,0) rectangle (\linewidth,-0.2\TPVertModule);
\draw[line width=2pt,color=ksline4] (0\linewidth,-0.2\TPVertModule) -- (1\linewidth,-0\TPVertModule);
\end{tikzpicture}
\end{textblock}

  % text box (free-positioned)
\begin{textblock}{0.05}(0.33,0.33)
{\raggedright\fontsize{18}{22}\selectfont \textcolor[HTML]{C8102E}{$r$}\par}
\end{textblock}

  % text box (free-positioned)
\begin{textblock}{0.05}(0.3,0.44)
{\raggedright\fontsize{16}{19}\selectfont $\theta$\par}
\end{textblock}

  % text box (free-positioned)
\begin{textblock}{0.44}(0.52,0.26)
{\raggedright\fontsize{16}{19}\selectfont \begin{itemize}
\item The integrand depends on $r$ only
\item Rings of radius $r$ have length $2\pi r$
\item That extra $r$ makes it integrable
\end{itemize}\par}
\end{textblock}
\end{frame}

% ======================================================================
% Slide 5 - Open questions
% ======================================================================
\begin{frame}[t]
  \frametitle{Open questions}

  % text box (free-positioned)
\begin{textblock}{0.88}(0.06,0.24)
{\raggedright\fontsize{19}{23}\selectfont \begin{enumerate}
\item A proof without polar coordinates?
\item Which other integrals yield to squaring?
\item Higher dimensions: $\int_{\mathbb{R}^n} e^{-|x|^2}\,dx = \pi^{n/2}$
\end{enumerate}\par}
\end{textblock}
\end{frame}

\end{document}
